\section{Structured cone systems and scalar estimation}\label{sec:structured}
A dense cone Hessian need not make its inverse forms hard to estimate. We
give two classical algorithms for structured Newton systems. The first
assumes a sparse positive definite base plus a few low-rank corrections,
with sampling access to the \emph{source} vectors of the corrections. It
returns a scalar estimate or a solution sampler without writing out a
dense Newton direction. The second uses an explicit factorization and
returns the full direction. They need different structural parameters.

The low-rank cone identities and their use in sparse factorizations are
established techniques \citep{AlizadehGoldfarb2003,GoldfarbScheinberg2005,ChenGoulart2025},
as is matrix arithmetic under sampling-and-query access with
dimension-independent cost \citep{ChiaEtAl2022}. The result below combines
a full-rank sparse base with possibly indefinite corrections. It needs its
own access model and error analysis because sampling from a nearly zero
transformed correction column can be arbitrarily expensive; the algorithm
never normalizes such a column.

\subsection{A sparse base with sampled low-rank corrections}
We work over $\R$ in this section. Suppose that $C\in\R^{m\times n}$ has
row and column sparsity at most $d$ and that public parameters give the
spectral bounds
\begin{equation}\label{eq:structured-model}
 S=CC^T,\quad \kappa^{-1}I\preceq S\preceq I,\qquad
 N=S+CXLX^TC^T,\quad aS\preceq N\preceq\beta S.
\end{equation}
Here $0<a\leq1\leq\beta$, $X\in\R^{n\times r}$ satisfies $\norm X\leq1$,
and the explicitly supplied $r\times r$ symmetric matrix $L$ obeys
$\norm L\leq\ell$, where $\ell\geq1$; $L$ need be neither invertible nor
semidefinite. Each column of $X$ is given as $X_j=T_jv_j$, where
$\norm{v_j}=1$, $\SQ(v_j)$ is supplied, and $T_j$ is a known contraction
with at most one entry per row and column, such as a diagonal map or a
coordinate embedding. Each $T_j$ has constant-cost sparse row and column
access as in Definition~\ref{def:access}, including reports of absent
entries. The norms of $X_j$ and $CX_j$ are not supplied. For the nonzero
right-hand side $b$, $\SQ(b)$ is supplied; by scaling, we assume
$\norm b=1$ throughout the proof.

\begin{theorem}[Sparse base and low-rank corrections]\label{thm:structured-sq}
Assume \eqref{eq:structured-model}, row and column sparse access to $C$,
and the vector access above. Then a randomized setup has an event of
probability at least $1-\zeta$, the good setup event, on which the
coordinate oracle it builds answers queries to one fixed vector
$\widetilde y$ with
\[
 \norm{\widetilde y-N^{-1}b}\leq\epsilon\norm{N^{-1}b},
\]
and on this event each requested draw returns either a failure flag or a
sample with law exactly $p_{\widetilde y}$ conditional on no flag, with
flag probability at most a prescribed $\delta\in(0,1/2)$.
Setup, coordinate queries, and draws have bounded cost even for bad
setup outcomes. With
$R_*=2+(r+1)\kappa\beta\ell/a$, each is at most
\begin{equation}\label{eq:structured-cost}
 (d+1)^{O(\sqrt\kappa\log(R_*/\epsilon))}
 \poly(R_*,\epsilon^{-1},\log(1/\zeta),\log(1/\delta)).
\end{equation}
No oracle for the exact norm of the output is provided. For
$r=0$ it reduces to Theorem~\ref{thm:solution-sampling} with a cap on the
number of rejection trials.
\end{theorem}
\begin{proof}
Put $D_*=1+\ell/a$ and $U_*=D_*\ell$. The matrices
\[
 K=X^TC^TS^{-1}CX,\quad h=X^TC^TS^{-1}b,\quad J=(I+LK)^{-1}
\]
obey $0\preceq K\preceq I$ and $\norm h\leq\sqrt\kappa$, since
$C^TS^{-1}C$ is an orthogonal projection. The identity
\begin{equation}\label{eq:structured-core}
 J=I-LX^TC^TN^{-1}CX,\qquad \norm J\leq D_*
\end{equation}
follows by multiplying by $I+LK$; the norm estimate uses
$N^{-1}\preceq a^{-1}S^{-1}$. Thus Woodbury's identity holds without
$L^{-1}$ and gives
\[
 z=JLh,\quad \norm z\leq U_*\sqrt\kappa,\qquad
 y=N^{-1}b=S^{-1}b-S^{-1}CXz.
\]

Use Lemma~\ref{lem:residual} to obtain a polynomial $P=P(S)$ with
$\norm{I-SP}\leq\eta\leq1$. Functional calculus gives
\begin{equation}\label{eq:structured-polynomial-norms}
 \norm P\leq2\kappa,\quad \norm{PC}\leq2\sqrt\kappa,
 \quad \norm{C^TPC}\leq2.
\end{equation}
Writing $C=S^{1/2}V$ with $VV^T=I$ proves all three.
Let $K_P=X^TC^TPCX$ and $h_P=X^TC^TPb$. Their biases satisfy
$\norm{K_P-K}\leq\eta$ and
$\norm{h_P-h}\leq\eta\sqrt\kappa$.

Each entry of $K_P$ is estimated by sampling from $v_i$ and evaluating
$T_i^TC^TPC T_jv_j$ locally. If $I\sim p_{v_i}$, the
variable $(T_i^TC^TPC T_jv_j)_I/(v_i)_I$ has the desired expectation and
second moment at most $4$; zero sampled denominators have probability zero.
For $(h_P)_i$ use $T_i^TC^TPb$ instead; its second moment is at most
$4\kappa$. Median-of-means estimates with entry errors
$\rho/(2r)$ and $\nu/(2\sqrt r)$, respectively, give
\[
 \norm{\widehat K-K}\leq\rho,\qquad
 \norm{\widehat h-h}\leq\nu
\]
with joint probability $1-\zeta$, provided
$\eta\leq\rho/2$ and $\eta\sqrt\kappa\leq\nu/2$.
One may symmetrize $\widehat K$ without increasing its error. The total
number of local evaluations is
\begin{equation}\label{eq:structured-estimates}
 O\left[\left(\frac{r^4}{\rho^2}
       +\frac{r^2\kappa}{\nu^2}\right)
       \log\frac{r+1}{\zeta}\right].
\end{equation}
All estimates sample from the supplied vectors, not from $CX_j$.

If $U_*\rho\leq1/2$, perturbing the $r\times r$ inverse gives
$\norm{(I+L\widehat K)^{-1}}\leq2D_*$ and, for
$\widehat z=(I+L\widehat K)^{-1}L\widehat h$,
\[
 \norm{\widehat z-z}\leq
 2U_*(\nu+\rho U_*\sqrt\kappa).
\]
The fixed output $\widetilde y=Pb-PCX\widehat z$ consequently satisfies
\begin{equation}\label{eq:structured-error}
 \norm{\widetilde y-y}\leq
 \eta\kappa(1+U_*)+4U_*\sqrt\kappa\,\nu
                   +4U_*^2\kappa\rho.
\end{equation}
Here the polynomial bias in $S^{-1}b$ is at most
$\eta\kappa$, that in $S^{-1}CXz$ is at most
$\eta\sqrt\kappa\norm z$, and the remaining term is bounded using
\eqref{eq:structured-polynomial-norms}. The explicit choices
\begin{align*}
 \rho&=\frac{\epsilon}{48\beta U_*^2\kappa},&
 \nu&=\frac{\epsilon}{48\beta U_*\sqrt\kappa},\\
 \eta&=\min\left\{\frac{\epsilon}{12\beta\kappa(1+U_*)},
                         \frac\rho2,\frac\nu{2\sqrt\kappa}\right\}
\end{align*}
make \eqref{eq:structured-error} at most $\epsilon/(4\beta)$.
Since $N\preceq\beta I$, $\norm y\geq1/\beta$, which proves the
stated relative error with room to spare. On the good event,
$\norm{\widehat z}\leq4U_*\sqrt\kappa$.

It remains to sample without dividing by small component norms. The raw
trial in Lemma~\ref{lem:rejection} also applies to a rectangular or singular
transform $T$: if its rows have at most $W$ entries and $B_T\geq\norm T$,
its output probability at coordinate $i$ in one trial is
\begin{equation}\label{eq:raw-rectangular}
                  \frac{|(Tv)_i|^2}{WB_T^2\norm v^2}.
\end{equation}
Its proof uses only column-norm upper bounds and the Cauchy--Schwarz
inequality on rows, so no lower singular-value bound is needed.
Choose one $W=(d+1)^{O(\sqrt\kappa\log(R_*/\epsilon))}$ bounding
the row and column support sizes of all transforms used and the cost of
computing their entries; sparse walks through $CC^T$ and the one-entry
maps $T_j$ give such a bound. Put $w_0=Pb$, $w_j=PC T_jv_j$, $c_0=1$ and
$c_j=-\widehat z_j$. Take the raw-trial denominators $D_0=4W\kappa^2$
and $D_j=4W\kappa$ for $j\geq1$, which equal $WB_T^2\norm v^2$ for the
bounds $B_T$ in \eqref{eq:structured-polynomial-norms}, and let
$Z=\sum_{j=0}^r c_j^2D_j$. Select $j$ with probability $c_j^2D_j/Z$
and perform one raw trial for $w_j$. If it returns $i$, accept again with
probability
\[
 \frac{|\sum_{j=0}^r c_jw_{j,i}|^2}
      {(r+1)\sum_{j=0}^r |c_jw_{j,i}|^2}.
\]
The denominator is positive whenever this step is reached. Summing over
$j$ shows that one combined trial outputs $i$ with probability
$|\widetilde y_i|^2/((r+1)Z)$. On the good setup event,
\[
 Z\leq4W\kappa^2(1+16U_*^2),\qquad
 \norm{\widetilde y}\geq\frac1{2\beta}.
\]
Thus the total acceptance probability is at least
$\pi_*=[16(r+1)W\beta^2\kappa^2(1+16U_*^2)]^{-1}$.

To bound the runtime on every outcome, setup returns a failure flag
if $I+L\widehat K$ is singular, if the Frobenius norm of its inverse
exceeds $2\sqrt r D_*$, or if $\norm{\widehat z}>4U_*\sqrt\kappa$. These
checks never reject a good setup and need only dense $r\times r$
arithmetic. Every draw is capped at
$\lceil\pi_*^{-1}\log(1/\delta)\rceil$ trials and returns a flag if all
fail. Since the trials are i.i.d., truncation does not change the law
conditional on acceptance. On a bad setup the loop still terminates, even if
$\widetilde y=0$. Each trial costs $O((r+1)W)$;
\eqref{eq:structured-estimates}, the $r\times r$ solve, and these caps
prove \eqref{eq:structured-cost}. For nonunit $b$, scale the fixed vector
back by the known norm of $b$; its sampling law is unchanged.
\end{proof}

\begin{corollary}[Structured positive inverse forms]\label{cor:structured-scalar}
Under the same hypotheses, $q=b^TN^{-1}b$ can be estimated to relative
error $\epsilon$ with failure probability at most $\zeta$, at a cost of
the form \eqref{eq:structured-cost}, using only setup and coordinate
queries and hence no per-draw failure parameter $\delta$.
\end{corollary}
\begin{proof}
Normalize $b$ to unit norm. Then $q\geq1/\beta$ and
$\norm{N^{-1}b}\leq\kappa/a$. Run the preceding setup with vector
error $\epsilon a/(4\beta\kappa)$ and failure budget $\zeta/2$.
On the good setup event, $|b^T\widetilde y-q|\leq\epsilon/(4\beta)$. Independently sample
$I\sim p_b$ and average $\widetilde y_I/b_I$ by median of means.
Conditionally on a good setup, its second moment is at most
$\norm{\widetilde y}^2\leq4\kappa^2/a^2$.
Thus $O(\beta^2\kappa^2a^{-2}\epsilon^{-2}\log(1/\zeta))$
coordinate evaluations achieve additional error $\epsilon/(2\beta)$.
A union bound proves the claim. The argument uses neither positivity of
the approximate linear map nor a promise that its individual correction
terms do not cancel.
\end{proof}

\subsection{Lorentz and generalized-power barriers}
At an interior Lorentz point $(t,z)$ with $t>\norm z$, let
$H=\nabla^2 F_L(t,z)$ for $F_L(t,z)=-\log(t^2-\norm z^2)$.
Put $r_z=\norm z$, $s=t^2-r_z^2>0$,
$d_x=s/2$, and $\chi=(t+r_z)/(t-r_z)$. If $r_z>0$, put
$u_\pm=2^{-1/2}(1,\pm z/r_z)$. Differentiation and inversion give
\begin{equation}\label{eq:lorentz-structured}
 H^{-1}=d_x I+r_z(t+r_z)u_+u_+^T-r_z(t-r_z)u_-u_-^T,
 \quad \chi^{-1}d_xI\preceq H^{-1}\preceq\chi d_xI.
\end{equation}
Indeed, the eigenvalues are $(t+r_z)^2/2$, $(t-r_z)^2/2$ and $s/2$
on the two radial eigendirections and their orthogonal complement. At
$z=0$ the correction vanishes. For $L_c$ blocks, let
$\mathcal D=\bigoplus d_xI$, $C=A\mathcal D^{1/2}$ and $S=CC^T$.
Form $X$ from the orthonormal columns $(1,0)$ and $(0,z/r_z)$ in each block.
The matrix $L$ in \eqref{eq:structured-model}, defined by
$\mathcal D^{-1/2}H^{-1}\mathcal D^{-1/2}=I+XLX^T$, is explicit, block
diagonal, and has norm at most $\chi_{\max}-1$, where $\chi_{\max}$ is the
largest $\chi$ over the blocks. Thus $r\leq2L_c$,
$a=\chi_{\max}^{-1}$ and $\beta=\chi_{\max}$ suffice in
\eqref{eq:structured-model}. Supplied block norms and $\SQ(z)$ give all
required source access. A known positive scaling of $S$ normalizes its
spectral upper bound to one and scales $N$ and the desired output
accordingly. Hence Theorem~\ref{thm:structured-sq} and
Corollary~\ref{cor:structured-scalar} apply to Lorentz cones.

For generalized-power cones, let $x_i>0$, $\alpha_i>0$,
$\sum_{i=1}^{m_x}\alpha_i=1$, and
\[
 p=\prod_i x_i^{2\alpha_i},\quad s=p-\norm z^2>0,
 \quad \lambda=p/s,\quad 1\leq\lambda\leq\Lambda.
\]
The barrier
\begin{equation}\label{eq:power-barrier}
 F(x,z)=-\log s-\sum_i(1-\alpha_i)\log x_i
\end{equation}
has parameter $m_x+1$ \citep{RoyXiao2022}. Its use and low-rank Hessian
structure are prior work \citep{ChenGoulart2025};
closed-form inverse-Hessian operators are given by
\citet{KapelevichAndersenVielma2024}.
Set $k_i=1+(2\lambda-1)\alpha_i$ and
$D_0=\diag(k_i/x_i^2,(2/s)I)$. With $c=\norm\alpha^2$,
define two columns with disjoint supports:
\[
 X_1=(\alpha_i/(\sqrt c\sqrt{k_i}))_i\oplus0,
 \qquad X_2=0\oplus z/\norm z.
\]
At $z=0$, omit the correction: $\lambda=1$ and $H=D_0$.
Otherwise direct differentiation of \eqref{eq:power-barrier} gives
\begin{equation}\label{eq:power-normalized}
 H=D_0^{1/2}(I+XBX^T)D_0^{1/2},\qquad
 B=\begin{pmatrix}
 4\lambda(\lambda-1)c&-4\lambda\sqrt{c(\lambda-1)/2}\\
 -4\lambda\sqrt{c(\lambda-1)/2}&2(\lambda-1)
 \end{pmatrix}.
\end{equation}
Before normalization, the correction is
$4\lambda(\lambda-1)aa^T-(4\lambda/s)(az^T+za^T)+(4/s^2)zz^T$,
where $a_i=\alpha_i/x_i$ and the $x$ and $z$ groups are embedded
separately; this confirms the signs in $B$ and the diagonal base $D_0$.

\begin{lemma}[Power-cone spectral comparison and estimated coefficient matrix]\label{lem:power-core}
The inverse Hessian satisfies
\begin{equation}\label{eq:power-comparison}
 (4\lambda)^{-1}D_0^{-1}\preceq H^{-1}\preceq4\lambda D_0^{-1}.
\end{equation}
In $H^{-1}=D_0^{-1/2}(I+XL(t)X^T)D_0^{-1/2}$, the coefficient
matrix is symmetric and given by
\[
 t=\sum_i\frac{\alpha_i^2}{c k_i},\quad
 L(t)=-(I+B\diag(t,1))^{-1}B.
\]
For every $t$ in the interval
\begin{equation}\label{eq:power-clip}
 \mathcal I_\lambda=
 [1/(2\lambda),\ \min\{1,1/(2\lambda c)\}],
\end{equation}
this formula is defined, $\norm{L(t)}\leq10\Lambda^2$, and
$\norm{L(t')-L(t)}\leq100\Lambda^4|t'-t|$.
The true value of $t$ lies in this interval and can be estimated to error
$u$ with $O(u^{-2}\log(1/\zeta))$ samples from $p_\alpha$.
\end{lemma}
\begin{proof}
Since $1\leq k_i\leq2\lambda$ and $k_i\geq2\lambda\alpha_i$,
$t\geq1/(2\lambda)$ and $ct\leq1/(2\lambda)$, proving the interval
claim. In the orthonormal basis obtained by normalizing the two column
directions of $X$, the normalized
Hessian is
\[
 Q(t)=\begin{pmatrix}
 1+4\lambda(\lambda-1)ct&-4\lambda\sqrt{ct(\lambda-1)/2}\\
 -4\lambda\sqrt{ct(\lambda-1)/2}&2\lambda-1
 \end{pmatrix}.
\]
Its trace is at most $4\lambda$ and its determinant is
$2\lambda-1-4\lambda(\lambda-1)ct\geq1$ on
\eqref{eq:power-clip}. Its leading diagonal entry is positive, so both
eigenvalues lie in $[1/(4\lambda),4\lambda]$. On the orthogonal
complement the matrix is the identity. This proves \eqref{eq:power-comparison}.
For $G=\diag(t,1)$,
\[
 L(t)=G^{-1/2}(Q(t)^{-1}-I)G^{-1/2}.
\]
This identity proves symmetry and gives
$\norm{L(t)}\leq2\lambda(4\lambda+1)\leq10\Lambda^2$.
Differentiating the inverse formula gives
$L'(t)=L(t)\diag(1,0)L(t)$, hence
$\norm{L'(t)}\leq100\Lambda^4$. Finally, $t$ is the expectation of
$1/k_I\in[0,1]$ for $I\sim p_\alpha$. Hoeffding's inequality proves the
sample bound. Clipping an estimate to \eqref{eq:power-clip} never increases
its error. Clipping only to $[1/(2\Lambda),1]$ would not preserve the
determinant bound and is insufficient.
\end{proof}

\begin{corollary}[Cone Newton systems with supplied block data]\label{cor:cone-sq}
Let $A$ have full row rank and row and column sparsity $d$, with
constant-cost sparse row and column access, a public block partition, and
constant-cost conversion between block-local and global indices. For $L_c$
Lorentz blocks, assume $\SQ(z)$ is supplied separately for each block and
axial coordinates can be queried; for $L_c$ generalized-power blocks,
assume $\SQ(\alpha)$ and $\SQ(z)$ are supplied separately for each block,
coordinates of $x$ can be queried, and each block's barrier scalars $p$
and $s$ are supplied. Thus each block's norm is available, not only the
norm and sampling law of the concatenated vector. Radial-vector norm and
coordinate access are available also when $z=0$; sampling is requested
only for blocks with nonzero radial vector, and a zero block's radial
correction is omitted. Supply $\SQ(b)$ and spectral bounds normalizing
$S=A\mathcal D A^T$ to $[\kappa_S^{-1},1]$, where $\mathcal D$ is
$\bigoplus d_xI$ (Lorentz) or $\bigoplus D_0^{-1}$ (power). Then, for
$N=AH^{-1}A^T$, relative approximation of $N^{-1}b$, capped sampling from
it, and relative estimation of $b^TN^{-1}b$ each cost
\begin{equation}\label{eq:cone-cost}
 (d+1)^{O(\sqrt{\kappa_S}\log(2L_c\Gamma\kappa_S/\epsilon))}
 \poly(L_c,\Gamma,\kappa_S,\epsilon^{-1},
                         \log(1/\zeta),\log(1/\delta)),
\end{equation}
where $\Gamma=\max\chi$ for Lorentz cones and $\Gamma=\Lambda$ for
power cones. The setup and draws have the guarantees of
Theorem~\ref{thm:structured-sq}; scalar estimation does not use $\delta$.
\end{corollary}
\begin{proof}
The Lorentz case has already been reduced to
Theorem~\ref{thm:structured-sq}. For power cones take $C=A\mathcal D^{1/2}$.
The columns in \eqref{eq:power-normalized} are diagonal contractions of
$\alpha/\norm\alpha$ and embeddings of $z/\norm z$; their supports are
disjoint across and within blocks, which gives $\norm X\leq1$. They do not
require the exact norm of $X_1$. Lemma~\ref{lem:power-core} gives
$a=1/(4\Lambda)$, $\beta=4\Lambda$, $\ell=10\Lambda^2$ and $r\leq2L_c$.
Only $L$ requires approximation. Estimate each $t$ and clip as in
\eqref{eq:power-clip}, obtaining a fixed block diagonal $\widehat L$.
If $\norm{\widehat L-L}\leq v$, then
\[
 \norm{S^{-1/2}(\widehat N-N)S^{-1/2}}\leq v,
 \quad \widehat N=S+CX\widehat L X^TC^T.
\]
For $v\leq a/2$, therefore,
$(a/2)S\preceq\widehat N\preceq(\beta+a/2)S$. The resolvent identity
also gives, for every $b$,
\[
 \norm{\widehat N^{-1}b-N^{-1}b}
 \leq\norm{\widehat N^{-1}}\norm{\widehat N-N}\norm{N^{-1}b}
 \leq\frac{2\kappa_S v}{a}\norm{N^{-1}b}.
\]
Take $v\leq\min\{a/2,\epsilon a/(8\kappa_S)\}$, and run the theorem
on $\widehat N$ with tolerance $\epsilon/4$. Since $\epsilon\leq1/2$,
the resulting total relative error is less than $\epsilon$.
The estimates of $t$ use accuracy $v/(100\Lambda^4)$ and joint failure
budget $\zeta/2$; the remaining setup uses $\zeta/2$. The lemma bounds
$\norm{\widehat L}\leq10\Lambda^2$ even for inaccurate estimates.
These are polynomial costs, and the logarithm of every required tolerance
has the order asserted in \eqref{eq:cone-cost}. For the scalar estimate,
run this vector construction at tolerance
$\epsilon a/(4\beta\kappa_S)$ and apply the independent estimator from
Corollary~\ref{cor:structured-scalar}.
\end{proof}

Corollary~\ref{cor:cone-sq} concerns one linear system with given oracles.
For fixed $d,L_c,\Gamma,\kappa_S$ and accuracy, growing only the cone
dimensions does not force a query cost polynomial in them. The corollary
does not cover maintaining the supplied access to the iterate or forming
$b$, and does not show that the parameters stay bounded along a path. The
cyclic one-cone example in Section~\ref{sec:realizations} has a growing
sparse-base condition number and hence does not contradict it. Standard
three-dimensional cone blocks give bounded block sparsity directly, so
Section~\ref{sec:classical} already applies when the reduced system has
the stated spectral bounds.

\subsection{The power-cone scalar oracle cannot be omitted}
Corollary~\ref{cor:cone-sq} assumes that each power block's $p$ is
supplied. Without it, obtaining $p$ from the vector access can require
$\Omega(M)$ queries. Fix uniform weights on $2M$ positive coordinates
grouped into pairs. Encode a hidden bit $h_j$ by choosing pair values
$(a,b)$ or $(c,d)$, where $a^2+b^2=c^2+d^2$ and $ab\ne cd$; the norm is
then public. A squared-coordinate sample is simulated by choosing a
uniform pair and querying its bit; a coordinate query also needs one bit.
For the quantum lower bound we fix one quantum version of this access:
coherent value queries, and a preparation unitary that prepares a uniform
pair index, queries its hidden bit, applies the corresponding public
two-coordinate rotation, and uncomputes the bit. Since all pairs have
equal norms, this prepares the normalized input vector. Coherent value
queries, the preparation unitary and its inverse, and their controlled
versions all use $O(1)$ hidden-bit queries. We make no claim about other
unitary completions of classical $\SQ$ access. However,
\[
 p=\big[(ab)^{M-|h|}(cd)^{|h|}\big]^{1/M}
\]
determines $|h|$ exactly and hence its parity. By the classical and
quantum bounded-error lower bounds for parity, computing $p$ exactly
requires $\Omega(M)$ queries under either access.
Even constant relative error is hard when the logarithmic range of the
coordinates is unbounded. Take
\[
 (a,b)=(\sqrt{1-e^{-2\gamma M}},e^{-\gamma M}),\qquad
 (c,d)=(1/\sqrt2,1/\sqrt2),\qquad \gamma=4.
\]
Both squared pair norms equal one, and adjacent possible values of $p$
have ratio
\[
 [(1/2)/(e^{-4M}\sqrt{1-e^{-8M}})]^{1/M}>e^4/2>3.
\]
A relative error less than $1/2$ therefore still determines the integer
$|h|$ and its parity, with the same quantum simulation. This proves the
same linear lower bound in ideal real arithmetic, using a range of
$\log x_i$ that grows with $M$. If that range is known and bounded, one
can instead estimate $p$ through the weighted mean of $\log x_i$. Samples
proportional to $\alpha_i$ suffice, but $\SQ(\alpha)$ samples
proportional to $\alpha_i^2$; the two coincide for uniform weights. These
costs are part of acquiring the input to Corollary~\ref{cor:cone-sq}.

The aggregated logarithm cone shows the limits of low-rank algebra alone.
In the interior of this cone, where $s=v\sum_i\log(w_i/v)-u>0$ and
$v,w_i>0$, the Hessian of
$-\log s-\log v-\sum_i\log w_i$ is diagonal plus rank at most three.
To see this, let $h=(0,0,(1/w_i)_i)$, and use the positive diagonal
with entries $1$, $m/(vs)+1/v^2$, and $(1+v/s)/w_i^2$, where $m$
is the number of $w$ coordinates. The remaining terms are
$-e_ue_u^T+(\nabla s)(\nabla s)^T/s^2-(e_vh^T+he_v^T)/s$;
$\nabla s$ belongs to $\operatorname{span}\{e_u,e_v,h\}$.
By the Woodbury identity, the inverse Hessian differs from the inverse
diagonal by rank at most three. This gives no uniform spectral comparison
like Lemma~\ref{lem:power-core} and no bound on the cost of obtaining the
needed scalars, so we make no analogue of Corollary~\ref{cor:cone-sq}
for this cone.

\subsection{Eccentricity profiles with full output}\label{subsec:profile}
We now turn to full-output algorithms. The first allows many Lorentz
cones, a few with arbitrarily large eccentricity $\chi_j$ (defined before
\eqref{eq:lorentz-structured}). Assume
$A=[A_1\ \cdots\ A_{L_c}]$ has full row rank, and each
$(t_j,z_j)$ is interior, with $t_j>\norm{z_j}$ and
$H_j=\nabla^2[-\log(t_j^2-\norm{z_j}^2)]$. Write
$N=\sum_j A_jH_j^{-1}A_j^T$ and
$S=\sum_j d_jA_jA_j^T$, with the notation of
\eqref{eq:lorentz-structured}. For a threshold $\theta\geq1$, keep the
blocks $B_\theta=\{j:\chi_j>\theta\}$ exact, replace the others by their
scalar base, and put
\[
 P_\theta=\sum_{j\in B_\theta}A_jH_j^{-1}A_j^T
              +\sum_{j\notin B_\theta}d_jA_jA_j^T,
 \qquad b_\theta=|B_\theta|.
\]
\begin{proposition}[Profile preconditioner and energy certificate]
\label{prop:profile}
The matrix $P_\theta$ is positive definite and
\begin{equation}\label{eq:profile-order}
 \theta^{-1}P_\theta\preceq N\preceq\theta P_\theta,
 \qquad \rank(P_\theta-S)\leq2b_\theta.
\end{equation}
Preconditioned conjugate gradients reduce the $N$-energy error by a factor
$\eta$ in $O(\theta\log(2/\eta))$ iterations. If $q=b-N\widehat y$,
then
\begin{equation}\label{eq:profile-residual}
 \theta^{-1}q^TP_\theta^{-1}q
 \leq\norm{\widehat y-N^{-1}b}_N^2
 \leq\theta q^TP_\theta^{-1}q.
\end{equation}
\end{proposition}
\begin{proof}
For each block $j\notin B_\theta$,
$\theta^{-1}d_jI\preceq H_j^{-1}\preceq\theta d_jI$.
The contributions of blocks in $B_\theta$ are identical in the two matrices and positive
semidefinite, so summing proves \eqref{eq:profile-order}.
Equation~\eqref{eq:lorentz-structured} gives the rank bound.
The preconditioned spectrum lies in $[\theta^{-1},\theta]$.
The usual Chebyshev residual bound for conjugate gradients is
$2((\sqrt\kappa-1)/(\sqrt\kappa+1))^k$ in the energy norm;
$\kappa\leq\theta^2$ proves the iteration count. If $\theta=1$,
one iteration suffices. Inverting the Loewner bounds and using
$\norm{\widehat y-N^{-1}b}_N^2=q^TN^{-1}q$ proves the certificate.
\end{proof}

The certificate has an exact Newton interpretation. For
$H\Delta x+A^T\Delta y=r_d$ and $A\Delta x=r_p$, the reduced right-hand
side is $b=AH^{-1}r_d-r_p$. An approximate multiplier induces
$\widehat{\Delta x}=H^{-1}(r_d-A^T\widehat{\Delta y})$, and
\begin{equation}\label{eq:primal-energy}
 \norm{\widehat{\Delta x}-\Delta x}_H^2
       =\norm{\widehat{\Delta y}-\Delta y}_N^2.
\end{equation}
Thus \eqref{eq:profile-residual} bounds the primal error in the barrier
norm, without a Euclidean condition-number factor. It does not by itself
bound the full KKT residual or a Euclidean residual.

We next bound the arithmetic cost of an explicit implementation. Let $A$
have $m$ rows and $n$ columns, suppose a width-$\tau$ elimination order
for the row-intersection graph of $A$ is supplied, and write
$\bar\tau=\tau+1$. Each column support
is a clique of size at most $\bar\tau$, so assembling $S$ costs
$O(\bar\tau\nnz A)$; its positive definite Cholesky factor costs
$O(m\bar\tau^2)$ and occupies $O(m\bar\tau)$ words. After deleting
zero update columns, write $P_\theta=S+UJU^T$ with at most $2b_\theta$
columns and diagonal $J$ with entries $\pm1$. The matrix
$J+U^TS^{-1}U$, of order at most $2b_\theta$, is nonsingular: the
determinant identity and $P_\theta\succ0$ prove this even when the update
is indefinite. Forming $U$, solving for $S^{-1}U$, forming its Gram
matrix, and factoring $J+U^TS^{-1}U$ give conservative exact-arithmetic
costs for setup and for each application of $P_\theta^{-1}$:
\begin{align}
 C_{\rm setup}&=O\bigl(n+\bar\tau\nnz A+m\bar\tau^2
                  +mb_\theta(\bar\tau+b_\theta)+b_\theta^3\bigr),
                  \label{eq:profile-setup}\\
 C_{\rm apply}&=O\bigl(m(\bar\tau+b_\theta)+b_\theta^2\bigr).
                  \label{eq:profile-apply}
\end{align}
These counts store the transformed correction columns densely.
A matrix-free product by $N$ costs $O(\nnz A+n)$ by applying
\eqref{eq:lorentz-structured} blockwise. Hence the total solve cost is
\begin{equation}\label{eq:profile-cost}
 C_{\rm prof}=C_{\rm setup}
     +O\bigl(\theta\log(2/\eta)[\nnz A+n+C_{\rm apply}]\bigr).
\end{equation}
To choose $\theta$ adaptively, sort the eccentricities in decreasing order.
Keeping the largest $b$ and choosing $\theta=\chi_{(b+1)}$, with
$\chi_{(L_c+1)}=1$, gives \eqref{eq:profile-cost} with correction count
at most $b$. Sorting costs $O(L_c\log L_c)$; evaluating the displayed
bound for all $b$ then costs $O(L_c)$. Ties do not affect the argument.
This trades the number of corrections against the largest remaining
eccentricity, in the supplied coordinates.

The rank bound $2b_\theta$ is sharp in the worst case. Take $b$
three-dimensional Lorentz blocks, $A=I$, and choose
$t+r_z=\sqrt{2\chi}$, $t-r_z=\sqrt{2/\chi}$. Then $S=I$ and $N$ has
three orthogonal $b$-dimensional eigenspaces $E_+,E_0,E_-$ with eigenvalues
$\chi,1,\chi^{-1}$. If $P=I+R\succ0$, $R=R^T$, and $\rank R<2b$,
then
\begin{equation}\label{eq:profile-rank-lower}
                  \kappa(P^{-1/2}NP^{-1/2})\geq\chi.
\end{equation}
To prove this, let $p_+,p_-$ be the positive and negative inertia of $R$.
If $p_+<b$, some $v\in E_+$ has $v^TRv\leq0$, so its generalized
Rayleigh quotient is at least $\chi$. Since $\dim\ker R>b$,
$\ker R\cap E_+^\perp$ contains a nonzero vector whose quotient is at
most one. Otherwise $p_+\geq b$ implies $p_-<b$; a vector in $E_-$ has
quotient at most $\chi^{-1}$ and one in $\ker R\cap E_-^\perp$ has
quotient at least one. This proves \eqref{eq:profile-rank-lower}.
Thus, among base-plus-update preconditioners, guaranteeing condition
number at most $\theta^2$ for every profile with $b$ blocks of
eccentricity $\chi>\theta^2$ requires update rank at least $2b$. In the
same family with $A=I$ and unequal eccentricities, the stronger Loewner
bounds \eqref{eq:profile-order} require at least $b_\theta$ positive and
$b_\theta$ negative update directions, by restriction to the high and low
eigenspaces of the blocks with $\chi_j>\theta$, so this rank bound is
sharp too. These are linear-algebra
bounds within a specified class of preconditioners, not oracle lower
bounds for optimization.

Signed Woodbury arithmetic can be numerically unstable. An alternative
uses $P_\theta=S+UU^T-VV^T$. Blockwise,
$d_jI-r_j(t_j-r_j)u_{j,-}u_{j,-}^T\succ0$, hence $S-VV^T\succ0$.
The augmented matrix
\[
 \begin{pmatrix}S&V&U\\V^T&I&0\\U^T&0&-I\end{pmatrix}
\]
is symmetric quasidefinite: its principal block on the first two
groups is positive definite, and its last block is negative definite.
Adding the at most $2b_\theta$ vertices for the columns of $V$ and $U$ to
every bag of a width-$\tau$ tree decomposition for $S$ increases the width
by at most $2b_\theta$. Quasidefinite matrices admit pivot-free
$LDL^T$ factorization under every symmetric permutation
\citep{Vanderbei1995}. The supplied width-$\tau$ order therefore gives
setup cost $O((m+b_\theta)(\bar\tau+2b_\theta)^2)$ and application and
storage cost $O((m+b_\theta)(\bar\tau+2b_\theta))$, plus coefficient
formation. Either implementation may be used.
Neither this count nor \eqref{eq:profile-cost} is a uniform
bit-complexity bound. Product-form factorizations address numerical
cancellation in this setting \citep{GoldfarbScheinberg2005}; working
precision and residual checks must still meet the accuracy required by
the outer solve.

\subsection{Latent sparsity and explicit-output comparisons}
Dense cone blocks can also have sparse representations with auxiliary
variables. For one Lorentz block, let
$J_0=\diag(1,-I)$, $u=J_0(t,z)$ and $s=t^2-\norm z^2$. Then
\begin{equation}\label{eq:onehub}
 H=-2J_0/s+4uu^T/s^2=D_x+ww^T,
 \qquad w=2u/s.
\end{equation}
The diagonal $D_x$ is indefinite. Collect these diagonals over all blocks
into $D_x$ and the embedded columns $w$ into $W$. The equality-constrained
Newton system has the same primal and multiplier solution as
\begin{equation}\label{eq:latent-augmentation}
 \begin{pmatrix}D_x&W&A^T\\W^T&-I&0\\A&0&0\end{pmatrix}
 \begin{pmatrix}\Delta x\\h\\\Delta y\end{pmatrix}
 =\begin{pmatrix}r_d\\0\\r_p\end{pmatrix}.
\end{equation}
We call the auxiliary variables $h$ hub variables, one per cone;
eliminating them gives $H$ exactly. Full row rank of $A$
and $H\succ0$ make both systems nonsingular. Define a graph with $n$
coordinate vertices, $m$ equality-row vertices, and two hub vertices per
cone (the two-hub factorization below uses both). Join row $k$ to
coordinate $i$ whenever $A_{ki}\ne0$, and join both hubs of a cone to all
coordinates of that cone. The matrix in
\eqref{eq:latent-augmentation} uses a subgraph of this fixed graph.

\begin{proposition}[Supplied latent width]\label{prop:latent-width}
If a width-$\tau_L$ tree decomposition of this graph, with
$O((\tau_L+1)(n+m+L_c))$ total bag entries, is supplied, then one Newton
system can be solved exactly in
\[
                 O((n+m+L_c)(\tau_L+1)^2)
\]
field operations, after $O(n+\nnz A)$ coefficient formation. This
concerns one right-hand side and asserts neither a reusable symmetric
factorization nor numerical stability.
\end{proposition}
\begin{proof}
Use \eqref{eq:onehub} to form \eqref{eq:latent-augmentation} without
square roots. Replace each vertex in each bag by its row and column
copies. Every nonzero of the matrix is now covered in its row--column
bipartite graph, both copies inherit a connected subtree, and the new
width is at most $2\tau_L+1$. Write $N_{\rm aug}=n+m+L_c$ for the
augmented matrix order and $k=O(\tau_L+1)$ for the copied decomposition
width plus one. The supplied bound on total bag entries implies
$|\mathcal T|=O(kN_{\rm aug})$ decomposition nodes.
The preprocessing in Section 3 of \citet{FurerHoppenTrevisan2025}
converts this compact decomposition to a nice decomposition with
$O(N_{\rm aug})$ nodes in
$O(k(|\mathcal T|+N_{\rm aug}))=O(k^2N_{\rm aug})$ time.
Their Corollary 3 then solves a consistent rectangular
system in $O((\hbox{rows}+\hbox{columns})(\hbox{width}+1)^2)$ field
operations. Our system is square and nonsingular; applying that corollary
proves the bound. Its arbitrary pivoting handles exact cancellations,
so no zero-free prescribed pivot order is assumed.
\end{proof}

For a reusable symmetric factorization, there is a two-hub alternative.
With $r_z>0$, define
\[
 D=2I/s,\quad
 u_+=\frac{\sqrt{2r_z(t+r_z)}}s(1,-z/r_z),\quad
 u_-=\frac{\sqrt{2r_z(t-r_z)}}s(1,z/r_z).
\]
Then $H=D+u_+u_+^T-u_-u_-^T$ and
$u_-^TD^{-1}u_-=2r_z/(t+r_z)<1$, so $D-u_-u_-^T\succ0$.
Zero radial vectors give zero hub columns. For the product barrier let
$U,V$ collect the positive and negative columns. Replacing the dual zero
block by $-\rho I$, $\rho>0$, gives the expanded matrix
\[
 \begin{pmatrix}
 D&V&U&A^T\\V^T&I&0&0\\U^T&0&-I&0\\A&0&0&-\rho I
 \end{pmatrix}.
\]
It is quasidefinite with positive group $(x,V)$ and negative group $(U,y)$.
A width-$\tau_L$ order gives factorization cost
$O((n+m+L_c)(\tau_L+1)^2)$, and storage and subsequent solve cost
$O((n+m+L_c)(\tau_L+1))$. Forming these coefficients uses norms and square
roots; these counts are in real arithmetic. If $K_0$ is the original KKT
matrix, $w_0=K_0^{-1}g\ne0$, and $\rho\norm{K_0^{-1}}<1$, the Neumann
bound gives
\[
 \frac{\norm{w_\rho-w_0}}{\norm{w_0}}
 \leq\frac{\rho\norm{K_0^{-1}}}{1-\rho\norm{K_0^{-1}}}.
\]
For $g=0$ both solutions are zero. A public bound on
$\norm{K_0^{-1}}$ thus controls the regularization bias.
Floating-point stability or successful refinement requires additional
scaling and backward-error hypotheses.

Two hubs per block are necessary if the base is a positive diagonal and
each hub adds a signed rank-one term on a generic block. Suppose all
coordinates of $u=J_0(t,z)$ are nonzero and the dimension is at least
three. Then a representation
$H=D'+\sigma vv^T$ with $D'\succ0$ diagonal would require
$\sigma(v_i/u_i)(v_j/u_j)=4/s^2$ for every $i\ne j$.
Three indices force all ratios to be equal and $\sigma=1$, leaving
$D'=-2J_0/s$, whose time diagonal is negative. This argument does not
apply to more general sparse bases.

A concrete graph class has constant latent width. Let the coarse bipartite
graph join each cone to the equality rows that use it, and assume that
each scalar column of $A$ occurs in at most one row. If the coarse graph
has width $w$, replace each cone vertex in each bag by its two hubs. For
a coordinate used by row $k$, attach a leaf bag $\{u_j,v_j,x_i,y_k\}$ to
a bag containing the corresponding coarse edge; unused coordinates need
only the first three vertices. Edge coverage and connectedness of vertex
bags are preserved. Consequently
\[
                    \tau_L\leq\max\{2w+1,3\}.
\]
Coarse forests have $\tau_L\leq3$, independently of block dimensions and
number of cones, although the materialized Hessians may contain
arbitrarily large cliques. The smaller of \eqref{eq:profile-cost} and the
bound of Proposition~\ref{prop:latent-width} is a valid exact-arithmetic
cost for a classical full output. The profile method tolerates large
latent width when few blocks are highly eccentric; the latent-width method
tolerates many eccentric blocks.

These comparisons concern the cost of one Newton solve, not whether a
different solver produces the same iterates. Replacing the solver in an
interior-point method that tolerates inexact directions requires the
structural and access assumptions, with uniform parameter bounds, at
every iterate the method may visit,
and a classical direction that meets the method's residual requirement.
For example, even without cone structure, for a consistent
full-column-rank system $Kx=f$ of total row and column dimension $D_K$,
with $M$ nonzeros and condition number $\kappa_K$, CGLS started from zero
produces iterates with
\[
 \norm{Kx_k-f}\leq
 2\left(\frac{\kappa_K-1}{\kappa_K+1}\right)^k\norm f.
\]
This follows by applying the CG energy bound to $K^TK$, whose condition number
is $\kappa_K^2$, and identifying that energy norm with the residual norm.
It gives the familiar exact-arithmetic cost
$O((M+D_K)\kappa_K\log(2/\eta))$ for relative residual $\eta$; this is not
a new result. For a positive definite system of order $D$, sparse direct
elimination with a supplied width-$w$ order costs $O(D(w+1)^2)$. A
quantum solver that returns the same explicit direction should be
compared with these elementary classical costs.

In particular, suppose $n+m=D$, $\nnz A\leq D^{1+o(1)}$, and either
profile parameters $\bar\tau,\theta,b_\theta=D^{o(1)}$ or a latent width
$\tau_L=D^{o(1)}$ are supplied. Then each Newton solve takes
$D^{1+o(1)}$ arithmetic operations, up to logarithmic factors in the
accuracy. In the latent-width case the bound on $\nnz A$ is automatic:
width $\tau_L$ bounds the edges of the structural incidence graph, hence
$\nnz A$, by $O(D(\tau_L+1))$, since $L_c\leq n$. Writing out a full
direction already takes $\Omega(D)$ word writes. So for these classes,
with the same input access for both solvers and these assumptions
holding uniformly over all iterates, replacing only the linear solver in a method that
writes out every iterate cannot give a speedup polynomial in the
dimension. This does not apply to scalar, sample, or quantum-state
outputs. Nor does a dense assembled matrix or a large raw condition number
by itself put an instance outside either bound.
